\section{বিশ্বগ্রাম}

\begin{learningbox}
এই টপিক শেষে শিক্ষার্থী পারবে:
\begin{itemize}
\item বিশ্বগ্রামের ধারণা ব্যাখ্যা করতে।
\item বিশ্বগ্রাম প্রতিষ্ঠায় তথ্য ও যোগাযোগ প্রযুক্তির ভূমিকা বিশ্লেষণ করতে।
\item শিক্ষা, চিকিৎসা, ব্যবসা, গবেষণা ও কর্মসংস্থানে বিশ্বগ্রামের প্রভাব উদাহরণসহ লিখতে।
\item বিশ্বগ্রামের সুবিধা ও ঝুঁকি আলাদা করতে।
\item বোর্ড পরীক্ষার সংক্ষিপ্ত ও সৃজনশীল প্রশ্নে প্রাসঙ্গিক উদাহরণ ব্যবহার করতে।
\end{itemize}
\end{learningbox}

\subsection{বিশ্বগ্রামের ধারণা}

তথ্য ও যোগাযোগ প্রযুক্তির কারণে পৃথিবীর দূরত্ব শুধু ভৌগোলিক দূরত্বে সীমাবদ্ধ থাকলেও মানুষের যোগাযোগ, শিক্ষা, চিকিৎসা, ব্যবসা ও সংস্কৃতি বিনিময় অনেক দ্রুত হয়েছে। ঢাকা থেকে একজন শিক্ষার্থী বিদেশের বিশ্ববিদ্যালয়ের lecture দেখতে পারে, একজন রোগী দূরের বিশেষজ্ঞ ডাক্তারের পরামর্শ নিতে পারে, আবার একজন freelancer নিজ দেশে বসে বিদেশি প্রতিষ্ঠানের কাজ করতে পারে। এই একে অপরের সঙ্গে দ্রুত যুক্ত হয়ে যাওয়ার ধারণাকেই বিশ্বগ্রাম বলা হয়।

\begin{definitionbox}{সংজ্ঞা}
বিশ্বগ্রাম হলো তথ্য ও যোগাযোগ প্রযুক্তি নির্ভর এমন একটি সামাজিক ও যোগাযোগব্যবস্থা, যেখানে পৃথিবীর বিভিন্ন স্থানের মানুষ দূরত্ব সত্ত্বেও দ্রুত যোগাযোগ, তথ্য বিনিময়, সেবা গ্রহণ ও সহযোগিতার মাধ্যমে একটি একক সমাজের মতো কাজ করতে পারে।
\end{definitionbox}

বিশ্বগ্রাম ধারণাটি যোগাযোগ তত্ত্ববিদ মার্শাল ম্যাকলুহানের সঙ্গে বেশি সম্পর্কিত। তাঁর ধারণার মূল কথা ছিল, ইলেকট্রনিক যোগাযোগমাধ্যম মানুষের পারস্পরিক সম্পর্ককে এমনভাবে বদলে দিচ্ছে, যেন পৃথিবী একটি ছোট গ্রামের মতো হয়ে যায়। আধুনিক ইন্টারনেট, মোবাইল যোগাযোগ, সামাজিক যোগাযোগমাধ্যম, ভিডিও কনফারেন্সিং, cloud service এবং e-commerce এই ধারণাকে বাস্তব জীবনে আরও স্পষ্ট করেছে।

\begin{keypointbox}{বিশ্বগ্রাম প্রতিষ্ঠার প্রধান উপাদান}
\begin{itemize}
\item \textbf{যোগাযোগ অবকাঠামো}: ইন্টারনেট, মোবাইল নেটওয়ার্ক, অপটিক্যাল ফাইবার, স্যাটেলাইট।
\item \textbf{ডিভাইস}: কম্পিউটার, স্মার্টফোন, ট্যাব, স্মার্ট টিভি, IoT device।
\item \textbf{সফটওয়্যার ও প্ল্যাটফর্ম}: ব্রাউজার, ই-মেইল, video meeting app, LMS, e-commerce platform।
\item \textbf{ডেটা ও কনটেন্ট}: text, image, audio, video, database, digital document।
\item \textbf{দক্ষতা}: digital literacy, cyber safety, online communication skill।
\item \textbf{নীতি ও নিরাপত্তা}: privacy, copyright, data protection, responsible use।
\end{itemize}
\end{keypointbox}

\begin{examplebox}{বাংলাদেশ প্রেক্ষিত}
বাংলাদেশে mobile banking, online admission, e-ticketing, telemedicine, online class, freelancing marketplace এবং সরকারি digital service বিশ্বগ্রামের বাস্তব উদাহরণ। এগুলো ব্যবহার করে মানুষ ঘরে বসেই তথ্য, অর্থ, শিক্ষা ও সেবা পেতে পারে।
\end{examplebox}

\subsection{যোগাযোগ ও টেলিকনফারেন্সিং}

বিশ্বগ্রামের সবচেয়ে দৃশ্যমান দিক হলো দ্রুত যোগাযোগ। আগে কোনো খবর বিদেশে পাঠাতে চিঠি বা fax-এর উপর নির্ভর করতে হতো; এখন ই-মেইল, instant messaging, voice call, video call, forum এবং social media ব্যবহার করে কয়েক সেকেন্ডে যোগাযোগ করা যায়।

\begin{definitionbox}{টেলিকনফারেন্সিং}
টেলিকনফারেন্সিং হলো এমন যোগাযোগব্যবস্থা, যেখানে ভিন্ন স্থানে থাকা একাধিক ব্যক্তি audio, video বা data communication technology ব্যবহার করে একই সময়ে সভা, ক্লাস, প্রশিক্ষণ বা আলোচনা করতে পারে।
\end{definitionbox}

\begin{center}
\begin{tabular}{|p{0.28\textwidth}|p{0.58\textwidth}|}
\hline
\textbf{যোগাযোগ মাধ্যম} & \textbf{ব্যবহার} \\
\hline
ই-মেইল & formal document, assignment, notice, file sharing \\
\hline
ভিডিও কনফারেন্স & online class, meeting, interview, medical consultation \\
\hline
সামাজিক যোগাযোগমাধ্যম & announcement, public awareness, community support \\
\hline
Cloud document & একই document একাধিক ব্যক্তি একসঙ্গে edit করা \\
\hline
\end{tabular}
\end{center}

\begin{boardbox}{বোর্ড ফোকাস}
“টেলিকনফারেন্সিং কীভাবে বিশ্বগ্রাম ধারণাকে বাস্তবায়ন করে?”—এ ধরনের প্রশ্নে দূরবর্তী যোগাযোগ, সময় ও খরচ সাশ্রয়, online meeting/class এবং instant information exchange উল্লেখ করলে উত্তর শক্তিশালী হয়।
\end{boardbox}

\subsection{কর্মসংস্থান}

বিশ্বগ্রামের কারণে কাজের বাজার স্থানীয় সীমার বাইরে চলে গেছে। একজন দক্ষ ব্যক্তি নিজের শহর বা গ্রামে থেকেও foreign client-এর কাজ করতে পারে। Web development, graphic design, content writing, data entry, digital marketing, app development, online tutoring ইত্যাদি ক্ষেত্রে outsourcing ও freelancing গুরুত্বপূর্ণ ভূমিকা রাখছে।

\begin{keypointbox}{কর্মসংস্থানে পরিবর্তন}
\begin{itemize}
\item remote work এবং work-from-home সুযোগ তৈরি হয়েছে।
\item freelancing marketplace তরুণদের আন্তর্জাতিক কাজের সুযোগ দিচ্ছে।
\item online portfolio, LinkedIn, GitHub, personal website দক্ষতা প্রদর্শনের মাধ্যম হয়েছে।
\item digital skill না থাকলে অনেক কাজের বাজারে পিছিয়ে পড়ার ঝুঁকি থাকে।
\end{itemize}
\end{keypointbox}

\begin{mistakebox}{সাধারণ ভুল}
অনেকে freelancing ও outsourcing একই অর্থে ব্যবহার করে। \textbf{Outsourcing} হলো কোনো কাজ বাইরের ব্যক্তি/প্রতিষ্ঠানকে দিয়ে করানো। \textbf{Freelancing} হলো স্বাধীনভাবে project-based কাজ করা। একজন freelancer outsourcing কাজ করতে পারে, কিন্তু দুটি শব্দ এক নয়।
\end{mistakebox}

\subsection{শিক্ষা}

বিশ্বগ্রাম শিক্ষাকে classroom-এর বাইরে নিয়ে গেছে। শিক্ষার্থী এখন video lecture, e-book, digital library, simulation, quiz platform, online exam, learning management system এবং virtual lab ব্যবহার করতে পারে। শিক্ষকও live class, recorded class, online feedback, shared document এবং discussion forum ব্যবহার করে শিক্ষার্থীদের সহায়তা করতে পারেন।

\begin{examplebox}{ই-লার্নিং উদাহরণ}
একজন HSC শিক্ষার্থী যদি সংখ্যা পদ্ধতির conversion বুঝতে না পারে, সে online video দেখে, interactive quiz দিয়ে, তারপর teacher-এর কাছে video call-এ প্রশ্ন করতে পারে। একই topic শেখার জন্য বই, video, quiz ও practice একসঙ্গে ব্যবহৃত হলে শেখা বেশি কার্যকর হয়।
\end{examplebox}

\begin{keypointbox}{শিক্ষায় বিশ্বগ্রামের সুবিধা}
\begin{itemize}
\item দূরবর্তী শিক্ষার্থী ভালো resource পেতে পারে।
\item দুর্বল শিক্ষার্থী recorded class বারবার দেখতে পারে।
\item শিক্ষক দ্রুত assignment ও feedback দিতে পারেন।
\item global course, webinar ও online certification পাওয়া যায়।
\end{itemize}
\end{keypointbox}

\subsection{চিকিৎসা ও টেলিমেডিসিন}

চিকিৎসাক্ষেত্রে বিশ্বগ্রামের বড় অবদান হলো টেলিমেডিসিন। দূরবর্তী রোগী audio/video call, electronic prescription, diagnostic report sharing, remote monitoring এবং specialist consultation-এর মাধ্যমে সেবা পেতে পারে। বিশ্ব স্বাস্থ্য সংস্থার digital health guideline অনুযায়ী telemedicine-এর মূল ধারণা হলো দূরত্বে থাকা মানুষের কাছে ICT ব্যবহার করে স্বাস্থ্যসেবা পৌঁছে দেওয়া।

\begin{definitionbox}{টেলিমেডিসিন}
টেলিমেডিসিন হলো তথ্য ও যোগাযোগ প্রযুক্তি ব্যবহার করে দূরবর্তী রোগীকে চিকিৎসা পরামর্শ, পর্যবেক্ষণ, রিপোর্ট বিশ্লেষণ বা follow-up সেবা দেওয়ার পদ্ধতি।
\end{definitionbox}

\begin{examplebox}{বাস্তব উদাহরণ}
গ্রামের একজন রোগী উপজেলা হাসপাতালে পরীক্ষা করিয়ে রিপোর্টের ছবি specialist doctor-কে পাঠাল। ডাক্তার video call-এ রিপোর্ট দেখে পরামর্শ দিলেন। এতে রোগীর সময়, ভ্রমণ খরচ ও ঝুঁকি কমল।
\end{examplebox}

\subsection{গবেষণা}

গবেষণায় বিশ্বগ্রাম জ্ঞান বিনিময়কে দ্রুত করেছে। গবেষকরা digital journal, online database, preprint server, cloud storage, collaborative document, video seminar এবং data analysis software ব্যবহার করে বিভিন্ন দেশের মানুষের সঙ্গে কাজ করতে পারেন। ICT ছাড়া আধুনিক বিজ্ঞান, চিকিৎসা, কৃষি, জলবায়ু, মহাকাশ বা bioinformatics গবেষণা কল্পনা করা কঠিন।

\begin{keypointbox}{গবেষণায় ICT-এর ভূমিকা}
\begin{itemize}
\item গবেষণা প্রবন্ধ ও data দ্রুত পাওয়া যায়।
\item একাধিক দেশের গবেষক একই project-এ কাজ করতে পারেন।
\item cloud ও supercomputing বড় data বিশ্লেষণে সহায়তা করে।
\item simulation ব্যবহার করে ঝুঁকিপূর্ণ বা ব্যয়বহুল পরীক্ষা আগে model করা যায়।
\end{itemize}
\end{keypointbox}

\subsection{অফিস}

আধুনিক office ব্যবস্থায় e-mail, e-file, document management system, project management software, video meeting, biometric attendance, ERP এবং cloud storage ব্যবহৃত হয়। ফলে office paperless, faster এবং accountable হতে পারে।

\begin{examplebox}{ই-অফিস}
একটি প্রতিষ্ঠানে leave application, bill approval, notice, meeting minutes এবং report যদি online system-এ জমা ও approve হয়, তাহলে সময় কম লাগে এবং document হারানোর ঝুঁকি কমে।
\end{examplebox}

\subsection{বাসস্থান}

বিশ্বগ্রাম শুধু অফিস বা বিদ্যালয় নয়; বাসস্থানেও ICT যুক্ত হয়েছে। smart home device, online bill payment, food delivery, e-ticketing, online shopping, home security camera, remote work setup এবং digital entertainment ঘরোয়া জীবনের অংশ হয়ে উঠেছে।

\begin{mistakebox}{সতর্কতা}
Smart device ব্যবহারের সুবিধা আছে, কিন্তু দুর্বল password, untrusted app এবং public Wi-Fi ব্যবহার করলে ব্যক্তিগত তথ্য ঝুঁকিতে পড়তে পারে। নিরাপদ password, two-factor authentication এবং নিয়মিত software update জরুরি।
\end{mistakebox}

\subsection{ব্যবসা-বাণিজ্য ও ই-কমার্স}

বিশ্বগ্রাম ব্যবসার ধরন বদলে দিয়েছে। দোকান বা showroom ছাড়াও website, social media page, marketplace app এবং mobile payment ব্যবহার করে পণ্য বিক্রি করা যায়। ক্রেতা online catalogue দেখে order করতে পারে, review পড়তে পারে, payment করতে পারে এবং delivery tracking করতে পারে।

\begin{definitionbox}{ই-কমার্স}
ইন্টারনেট ও digital payment system ব্যবহার করে পণ্য বা সেবা কেনাবেচা করার প্রক্রিয়াকে ই-কমার্স বলা হয়।
\end{definitionbox}

\begin{center}
\begin{tabular}{|p{0.26\textwidth}|p{0.60\textwidth}|}
\hline
\textbf{ক্ষেত্র} & \textbf{বিশ্বগ্রামের প্রভাব} \\
\hline
ক্ষুদ্র ব্যবসা & social media ও marketplace-এ কম খরচে ক্রেতা পাওয়া \\
\hline
ব্যাংকিং & mobile banking, online transfer, bill payment \\
\hline
সরবরাহ ব্যবস্থা & inventory software, delivery tracking, online order \\
\hline
গ্রাহকসেবা & chatbot, review, return request, helpdesk \\
\hline
\end{tabular}
\end{center}

\subsection{সংবাদ}

Online news portal, live streaming, social media update এবং citizen journalism-এর মাধ্যমে খবর দ্রুত ছড়ায়। দুর্যোগ, পরীক্ষার notice, সরকারি announcement, traffic update বা emergency information দ্রুত মানুষের কাছে পৌঁছে যায়।

\begin{mistakebox}{ভুল তথ্য থেকে সাবধান}
দ্রুত সংবাদ পাওয়ার সুবিধা থাকলেও fake news ছড়ানোর ঝুঁকি আছে। কোনো খবর share করার আগে source, date, original statement এবং একাধিক reliable source মিলিয়ে দেখা উচিত।
\end{mistakebox}

\subsection{বিনোদন ও সামাজিক যোগাযোগ}

OTT platform, online game, music streaming, social media, short video platform এবং virtual community মানুষের বিনোদন ও সামাজিক সম্পর্কের ধরন বদলে দিয়েছে। মানুষ একই interest-এর community-তে যুক্ত হতে পারে, live event দেখতে পারে এবং নিজের creativity প্রকাশ করতে পারে।

তবে অতিরিক্ত screen time, cyberbullying, privacy leak, misinformation এবং attention loss-এর ঝুঁকিও আছে। তাই digital wellbeing ও responsible use শেখা জরুরি।

\subsection{সাংস্কৃতিক বিনিময়}

বিশ্বগ্রাম সংস্কৃতি বিনিময়কে সহজ করেছে। মানুষ online platform-এ ভাষা, গান, চলচ্চিত্র, উৎসব, পোশাক, খাবার, সাহিত্য ও লোকজ ঐতিহ্য সম্পর্কে জানতে পারে। বাংলাদেশের সংস্কৃতি, মুক্তিযুদ্ধের ইতিহাস, লোকসংগীত, পিঠা উৎসব, জামদানি, নকশিকাঁথা ইত্যাদিও digital archive, documentary, website ও social media-এর মাধ্যমে বিশ্বের কাছে পৌঁছাতে পারে।

\begin{recapbox}
\begin{itemize}
\item বিশ্বগ্রাম মানে ICT-এর মাধ্যমে পৃথিবীর মানুষ একক সমাজের মতো যুক্ত হওয়া।
\item যোগাযোগ, শিক্ষা, চিকিৎসা, গবেষণা, ব্যবসা ও সংস্কৃতি—সব জায়গায় এর প্রভাব আছে।
\item Internet, device, software, data, digital skill ও cyber safety বিশ্বগ্রামের ভিত্তি।
\item সুবিধার পাশাপাশি privacy, misinformation, digital divide ও cybercrime ঝুঁকি আছে।
\item পরীক্ষায় উদাহরণ দিলে বাংলাদেশ প্রেক্ষিত ব্যবহার করলে উত্তর বেশি শক্তিশালী হয়।
\end{itemize}
\end{recapbox}

\begin{boardbox}{পরীক্ষায় যেভাবে আসতে পারে}
\begin{enumerate}
\item বিশ্বগ্রাম বলতে কী বোঝায়?
\item বিশ্বগ্রাম প্রতিষ্ঠায় ICT-এর ভূমিকা ব্যাখ্যা কর।
\item “টেলিমেডিসিন বিশ্বগ্রামের একটি বাস্তব উদাহরণ”—ব্যাখ্যা কর।
\item বিশ্বগ্রাম কর্মসংস্থান ও শিক্ষায় কী পরিবর্তন এনেছে?
\item বিশ্বগ্রামের সুবিধা ও অসুবিধা তুলনা কর।
\end{enumerate}
\end{boardbox}

\begin{practicebox}
\textbf{MCQ}
\begin{enumerate}
\item বিশ্বগ্রাম ধারণার সঙ্গে সবচেয়ে বেশি সম্পর্কিত ব্যক্তি কে?\\
ক. বিল গেটস \quad খ. মার্শাল ম্যাকলুহান \quad গ. চার্লস ব্যাবেজ \quad ঘ. টিম বার্নার্স-লি
\item দূরবর্তী চিকিৎসাসেবা প্রদানের ICT-ভিত্তিক পদ্ধতি কোনটি?\\
ক. ই-কমার্স \quad খ. টেলিমেডিসিন \quad গ. ই-অফিস \quad ঘ. ব্লগিং
\end{enumerate}

\textbf{সংক্ষিপ্ত প্রশ্ন}
\begin{enumerate}
\item বিশ্বগ্রামের দুটি সুবিধা লেখ।
\item ই-কমার্স কীভাবে ক্ষুদ্র ব্যবসাকে সহায়তা করে?
\end{enumerate}

\textbf{সৃজনশীল অনুশীলন}
একজন শিক্ষার্থী গ্রামে বসে online class করে, mobile banking দিয়ে course fee দেয় এবং বিদেশি client-এর জন্য logo design করে আয় করে। উদ্দীপকের আলোকে বিশ্বগ্রামের তিনটি প্রভাব বিশ্লেষণ কর।
\end{practicebox}

\begin{sourcebox}
এই topic লেখায় local resource PDF, NCTB-style chapter outline, Marshall McLuhan-এর global village ধারণা, Wikipedia-এর Global Village overview, WHO digital health guideline-এর telemedicine অংশ এবং HSC ICT online topic index মিলিয়ে নিজের ভাষায় ব্যাখ্যা করা হয়েছে।
\end{sourcebox}
